
Soit @@M0@@ un @@M1@@-Module localement libre de rang fini @@M2@@. Il existe un torseur @@M3@@ sous @@M4@@ tel que @@M5@@, où @@M6@@ est considéré comme @@M7@@-Module par la représentation standard. On peut alors définir un @@M8@@-homomorphisme de @@M9@@ dans @@M10@@, en associant à toute représentation @@M11@@ de @@M12@@ dans un @@M13@@-module libre de rang @@M14@@ la classe du @@M15@@-Module localement libre
@@M16@@, où @@M17@@ est considéré comme @@M18@@-module par la représentation @@M19@@. Dans cet homomorphisme, la classe de @@M20@@ est donc l'image de la représentation standard. Pour vérifier une relation polynomiale entre les @@M21@@, il suffit donc que cette relation soit vérifiée par la représentation standard dans @@M22@@.

On peut faire un raisonnement analogue avec deux @@M23@@-Modules localement libres @@M24@@ de rangs @@M25@@ ; on peut alors trouver un torseur @@M26@@ sous @@M27@@ tel que @@M28@@, où @@M29@@ est considéré comme @@M30@@-Module en faisant opérer @@M31@@ par la représentation standard et @@M32@@ par la représentation triviale, et réciproquement pour @@M33@@. On définit ensuite un @@M34@@-homomorphisme de @@M35@@ dans @@M36@@ de la même façon que plus haut, et on se ramène ainsi à vérifier les relations entre les @@M37@@, @@M38@@ pour les éléments correspondants de @@M39@@.

On est donc ramené à montrer que pour tout schéma en groupes @@M40@@ sur @@M41@@, produit fini de schémas en groupes linéaires, le pré-@@M42@@-anneau @@M43@@ est un @@M44@@-anneau. Pour ce faire, on peut référer à [5], où @@M45@@ est calculé explicitement. Si l'on ne veut pas utiliser ce résultat, on peut raisonner de façon analogue à 3.2, en utilisant 1.14.

